\chapter{Ciąg dokładny chirurgii i zbiór struktur}
\label{ch:ciag-chirurgii}
\index{ciag dokladny chirurgii@Ciąg dokładny chirurgii}

\noindent\textbf{Pytanie przewodnie.}
Twierdzenie~\ref{thm:Wall-L-22} rozstrzyga, czy
\emph{jedno} normalne odwzorowanie można poprawić
do prostej równoważności homotopijnej.
Do klasyfikacji potrzebujemy jeszcze wiedzieć,
które odwzorowania normalne liczymy i kiedy
dwie otrzymane równoważności opisują tę samą
gładką rozmaitość. Odpowiedź łączy normalne
niezmienniki, grupy $L$ i zbiór struktur.
W całym rozdziale, o ile nie zaznaczono inaczej,
$X^n$ jest spójną, zamkniętą, gładką
rozmaitością wymiaru $n\geq6$.
Ta konwencja pozwala stosować gładką
wersję Twierdzenia~\ref{thm:s-kobordyzm-20}.

\section{Od normalnego odwzorowania do normalnego niezmiennika}
\label{sec:normalne-niezmienniki-23}

Definicja~\ref{def:normal-map-21} wiąże
odwzorowanie $f:M\to X$ z izomorfizmem
stabilnych wiązek normalnych nad pewną
wiązką $\xi\to X$. Dwie takie trójki
uważamy za tę samą \emph{daną normalną},
gdy łączy je normalny kobordyzm z
Definicji~\ref{def:normal-bordism-21}.
Nie żądamy jeszcze, aby $f$ było
równoważnością homotopijną.

\begin{definicja}[Zbiór normalnych niezmienników]
\index{zbior normalnych niezmiennikow@Zbiór normalnych niezmienników}
\label{def:normal-invariants-23}
Przez $\mathcal N_{\mathrm{DIFF}}(X)$ oznaczamy
zbiór klas normalnych odwzorowań stopnia jeden
$(f,b,\xi):M^n\to X^n$, gdzie wiązka
$\xi\to X$ może się zmieniać, względem normalnego
kobordyzmu. Punktem wyróżnionym jest
$(\id_X,\id_{\nu_X},\nu_X)$.
\end{definicja}

\begin{przyklad}[Identyczność]
\label{prz:normalny-punkt-23}
Identyczność $X\to X$ ma stopień jeden i
kanoniczny izomorfizm wiązek normalnych
przy wyborze $\xi=\nu_X$.
Stanowi zero zbioru normalnych niezmienników.
Z kolei odwzorowanie stopnia jeden bez
izomorfizmu~\eqref{eq:normal-map-21}
nie określa elementu $\mathcal N_{\mathrm{DIFF}}(X)$;
sam stopień nie pamięta obramowań
potrzebnych do chirurgii.
\end{przyklad}

Aby zapisać te dane homotopijnie, niech
$O=\operatorname*{colim}_{q}O(q)$, a $G$ oznacza
stabilną przestrzeń samorównoważności
homotopijnych sfer (składowe stopnia $\pm1$).
Odwzorowanie $O\to G$ zapomina liniową
strukturę obrotu sfery. Jego stabilny iloraz
homotopijny oznaczamy $G/O$.
W równoważnym opisie jest on włóknem
homotopijnym $BO\to BG$: mierzy wybór
stabilnej struktury wektorowej na zadanej
stabilnej sferycznej fibracji.
Na przykład dwie konstrukcje z tą samą
fibracją sferyczną mogą mieć różne
stabilne wiązki wektorowe; różnicę rejestruje
właśnie klasa w $G/O$.

\begin{twierdzenie}[Interpretacja normalnych niezmienników;
 wynik zewnętrzny]
\label{thm:normal-invariants-GO-23}
Dla zamkniętej gładkiej rozmaitości $X$
konstrukcja Pontriagina--Thoma daje
naturalną bijekcję zbiorów z punktami
wyróżnionymi
\[
 \mathcal N_{\mathrm{DIFF}}(X)\cong [X,G/O].
\]
Prawa strona oznacza klasy homotopii
odwzorowań do $G/O$; punkt stały odpowiada
identyczności $X\to X$.
\end{twierdzenie}

\noindent\textbf{Co stoi za bijekcją.}
Zanurzamy $M$ w wysokowymiarowej przestrzeni
euklidesowej i zwijamy dopełnienie jej
otoczenia rurkowego do punktu.
Izomorfizm $b$ wraz z warunkiem stopnia jeden
porównuje otrzymaną fibrację sferyczną z tą
nad $X$; wybór $\xi$ może się zmieniać.
Kobordyzm normalny staje się homotopią
wyborów podniesienia do $BO$.
W drugą stronę transwersalna reprezentacja
klasy $[X,G/O]$ daje podrozmaitość i
odwzorowanie stopnia jeden.
Twierdzenie transwersalności i stabilny
rachunek Pontriagina--Thoma są tutaj
wynikami zewnętrznymi; pełny dowód
identyfikacji znajduje się u
\href{https://webhomes.maths.ed.ac.uk/~v1ranick/books/scm.pdf}
{Walla, \emph{Surgery on Compact Manifolds},
lemat 10.6}.
Ten opis wyjaśnia intuicję, nie zastępuje
dowodu konstrukcji odwrotnej.

\section{Jak rozróżnić gładkie realizacje celu?}
\label{sec:zbior-struktur-23}

Wcześniej, w Rozdziale~\ref{ch:s-kobordyzm},
odróżniliśmy zwykłą od prostej równoważności.
To rozróżnienie zachowujemy w klasyfikacji.

\begin{definicja}[Prosty gładki zbiór struktur]
\index{prosty gladki zbior struktur@Prosty gładki zbiór struktur}
\label{def:structure-set-23}
$\mathcal S^s_{\mathrm{DIFF}}(X)$ jest zbiorem
klas par $(M,f)$, gdzie $M$ jest zamkniętą
gładką rozmaitością, a $f:M\to X$ prostą
równoważnością homotopijną. Pary $(M_0,f_0)$
i $(M_1,f_1)$ są równoważne, gdy istnieje
dyfeomorfizm $h:M_0\to M_1$ taki, że
$f_1\circ h\simeq f_0$. Punktem wyróżnionym
jest $(X,\id_X)$.
\end{definicja}

Intuicja: element $\mathcal S^s_{\mathrm{DIFF}}(X)$
to gładka realizacja ustalonego typu homotopijnego
razem z identyfikacją tego typu z $X$.

\begin{przyklad}[Ten sam typ homotopijny]
\label{prz:struktury-sfera-23}
Jeśli $\Sigma^n$ jest gładką sferą homotopijną,
każda wybrana prosta równoważność
$\Sigma\to S^n$ daje element
$\mathcal S^s_{\mathrm{DIFF}}(S^n)$.
Samo stwierdzenie $\Sigma\simeq S^n$
nie orzeka, czy ten element jest punktem
wyróżnionym: definicja wymaga dyfeomorfizmu.
Konkretne niezerowe przykłady zobaczymy
w Rozdziale~\ref{ch:sfery-Milnora}.
\end{przyklad}

Prosta równoważność $f$ ma kanoniczną klasę
normalną: przenosi stabilną normalną
fibrację sferyczną, a stabilną wiązkę
$\nu_M$ można przenieść na wiązkę $\xi$
nad $X$ za pomocą odwrotności homotopijnej $f$.
Tak otrzymana wiązka nie musi być
stabilnie izomorficzna z $\nu_X$.
Wybór wyznacza element $\mathcal N_{\mathrm{DIFF}}(X)$.
Otrzymujemy odwzorowanie
\[
 \eta:\mathcal S^s_{\mathrm{DIFF}}(X)
          \longrightarrow\mathcal N_{\mathrm{DIFF}}(X).
\]
To nie jest dowolne zapomnienie:
jeżeli dwie struktury mają to samo $\eta$,
istnieje między nimi normalny kobordyzm
nad $X\times I$. Pozostała przeszkoda
tego kobordyzmu leży w grupie $L_{n+1}$.

\section{Działanie grupy \texorpdfstring{$L$}{L} i ciąg dokładny}
\label{sec:exact-surgery-23}

Niech $\pi=\pi_1(X)$, a
$w:\pi\to\{\pm1\}$ będzie charakterem
orientacji. Twierdzenie~\ref{thm:Wall-L-22}
określa odwzorowanie
\[
 \theta:\mathcal N_{\mathrm{DIFF}}(X)
       \longrightarrow L_n^s(\Z[\pi],w).
\]
Zbiór $[X,G/O]$ ma naturalną strukturę grupy
abelowej, lecz $\theta$ na ogół nie jest
homomorfizmem tej grupy. W ciągu używamy go
jako odwzorowania zbiorów z punktami wyróżnionymi;
do pojęcia dokładności wystarcza przeciwobraz zera.

\begin{twierdzenie}[Ciąg dokładny chirurgii;
 wynik Walla]
\label{thm:exact-surgery-23}
Dla $X$ jak na początku rozdziału istnieje
ciąg zbiorów z punktami wyróżnionymi
\begin{equation}\label{eq:exact-surgery-23}
 [\Sigma X_+,G/O]\xrightarrow{\ \theta_{n+1}\ }
 L_{n+1}^s(\Z[\pi],w)
 \xrightarrow{\ \partial_{\!s}\ }
 \mathcal S^s_{\mathrm{DIFF}}(X)
 \xrightarrow{\ \eta\ }
 [X,G/O]
 \xrightarrow{\ \theta_n\ }
 L_n^s(\Z[\pi],w).
\end{equation}
Tutaj $X_+$ oznacza $X$ z dodanym punktem
bazowym, a $\Sigma$ zawieszenie zredukowane.
W odróżnieniu od $\theta_n$ lewa strzałka
$\theta_{n+1}$ jest homomorfizmem grup;
to osobny wniosek z konstrukcji na cylindrze
(Wall, stwierdzenie~10.7).
Grupa $L_{n+1}^s$ działa na zbiorze struktur;
$\partial_{\!s}$ wysyła jej element na wynik
działania na strukturze $(X,\id_X)$.
Dokładność oznacza:
\begin{enumerate}[label=(\roman*)]
\item obraz $\eta$ jest dokładnie
      $\theta_n^{-1}(0)$;
\item dwie struktury mają ten sam obraz
      przez $\eta$ wtedy i tylko wtedy,
      gdy leżą w jednej orbicie działania
      $L_{n+1}^s$;
\item stabilizator punktu wyróżnionego
      w $L_{n+1}^s$ jest obrazem
      $\theta_{n+1}$.
\end{enumerate}
\end{twierdzenie}

\begin{proof}[Wyjaśnienie kolejnych miejsc dokładności]
Najpierw weźmy normalną klasę $a$
z $\theta_n(a)=0$. Twierdzenie Walla
daje normalny kobordyzm od jej
reprezentanta do prostej równoważności.
Ta równoważność określa strukturę
z $\eta([M,f])=a$. Odwrotnie struktura
jest już prostą równoważnością,
więc jej przeszkoda znika.
To dowodzi (i).

Jeśli $\eta([M_0,f_0])=\eta([M_1,f_1])$,
normalny kobordyzm między dwiema
strukturami jest problemem chirurgicznym
nad $X\times I$. Jego przeszkoda ma
wymiar $n+1$ i należy do $L_{n+1}^s$.
Geometryczne twierdzenie realizacyjne
Walla pozwala zrealizować każdy taki
element przez odpowiedni kobordyzm;
operacja na górnym brzegu daje działanie.
Po zniwelowaniu przeszkody otrzymujemy prostą
równoważność kobordyzmu nad $X\times I$,
relatywnie do obu końców. Ponieważ oba końce
są już prostymi równoważnościami z $X$,
ich inkluzje do kobordyzmu są prostymi
równoważnościami; powstaje $s$-kobordyzm, a
Twierdzenie~\ref{thm:s-kobordyzm-20}
zamienia go w dyfeomorfizm końców.
W drugą stronę działanie zachowuje
normalną klasę, bo jego ślad jest
normalnym kobordyzmem. To daje (ii).

Jeśli element $L_{n+1}^s$ pochodzi
z normalnej klasy na cylindrze
$X\times I$ ustalonej na obu końcach,
jego działanie na $(X,\id_X)$
zwraca tę samą strukturę.
Twierdzenie o realizacji i
odwrotne rozpoznanie takiego
cylindra dają również przeciwną
implikację; jest to (iii).
Utożsamienie normalnych danych cylindra
z $[\Sigma X_+,G/O]$ wynika
z ilorazu $X\times I/(X\times\partial I)$.
Te głębokie kroki są odpowiednio
twierdzeniami 10.4, 10.5 i 10.8 w
\href{https://webhomes.maths.ed.ac.uk/~v1ranick/books/scm.pdf}
{książce Walla, §10}; przywołaliśmy
ich hipotezy i zastosowanie, a nie
powtórzyliśmy ich pełnego dowodu.
\end{proof}

\begin{figure}[htbp]
\centering
\begin{tikzpicture}[>=Latex,font=\small,
 every node/.style={align=center}]
 \node[draw=manifoldblue!75!black,fill=manifoldblue!8,
   rounded corners,minimum width=2.15cm,minimum height=.9cm]
   (l1) at (0,0) {$L_{n+1}^s$};
 \node[draw=manifoldblue!75!black,fill=manifoldblue!8,
   rounded corners,minimum width=2.15cm,minimum height=.9cm]
   (s) at (3.5,0) {$\mathcal S^s(X)$};
 \node[draw=manifoldblue!75!black,fill=manifoldblue!8,
   rounded corners,minimum width=2.15cm,minimum height=.9cm]
   (g) at (7.0,0) {$[X,G/O]$};
 \node[draw=manifoldblue!75!black,fill=manifoldblue!8,
   rounded corners,minimum width=2.15cm,minimum height=.9cm]
   (l2) at (10.5,0) {$L_n^s$};
 \draw[->,very thick,manifoldblue!70!black]
   (l1)--node[above] {działa}(s);
 \draw[->,very thick,manifoldblue!70!black]
   (s)--node[above] {$\eta$}(g);
 \draw[->,very thick,manifoldblue!70!black]
   (g)--node[above] {$\theta$}(l2);
 \node[font=\scriptsize,manifoldred!80!black]
   at (8.75,-1.0) {$\theta=0$: realizacja};
 \node[font=\scriptsize,manifoldgold!65!black]
   at (3.5,-1.0) {włókno $\eta$ jest orbitą};
\end{tikzpicture}
\caption{Środkowa część ciągu chirurgii.
Zbiór struktur nie jest zwykle grupą; dokładność
opisuje realizowalność normalnych danych oraz
orbity działania grupy $L_{n+1}^s$.}
\label{fig:exact-surgery-23}
\end{figure}

\section{Co ciąg mówi w praktyce?}
\label{sec:applications-surgery-23}

\begin{przyklad}[Cel po prostu spójny]
\label{prz:simply-connected-surgery-23}
Gdy $\pi=1$, obliczenie z
Sekcji~\ref{sec:obliczenia-LZ-22} daje
\[
 L_{4j}^s(\Z)\cong\Z,\qquad
 L_{4j+2}^s(\Z)\cong\Z/2,\qquad
 L_{2j+1}^s(\Z)=0.
\]
W wymiarze $4j$ pierwsza przeszkoda
normalnego odwzorowania jest sygnaturą
formy jądra podzieloną przez $8$.
Rozszczepienie z
Lematu~\ref{lem:split-normalne-21}
jest ortogonalne dla przecięć: dla
$y\in K=\ker f_*$ i $x\in H_{2j}(X;\Q)$
naturalność przecięć daje
$\lambda_M(f_!x,y)=\lambda_X(x,f_*y)=0$.
Dlatego
sygnatura jądra to
$\operatorname{sign}(M)-\operatorname{sign}(X)$.
W wymiarze $4j+2$ mierzy ją
niezmiennik Arfa formy kwadratowej.
W wymiarach nieparzystych grupa przeszkód
jest zerowa. W każdym z tych przypadków
ciąg nadal wymaga obliczenia $[X,G/O]$
oraz działania $L_{n+1}$, zanim
rozróżnimy wszystkie struktury.
\end{przyklad}

\begin{przyklad}[Dlaczego zerowa przeszkoda
nie oznacza jedynej struktury]
\label{prz:nonunique-structure-23}
Dla $\theta_n(a)=0$ istnieje
\emph{co najmniej jedna} prosta równoważność
o normalnej klasie $a$. Wszystkie
struktury z tą klasą tworzą jedną orbitę
$L_{n+1}^s$, która może zawierać
więcej niż jeden element.
To jest algebraiczne miejsce na różne
gładkie realizacje tego samego typu
homotopijnego. Sfery Milnora z
Rozdziału~\ref{ch:sfery-Milnora}
dadzą konkretny powód, dla którego
samej równoważności homotopijnej
nie wolno utożsamiać z dyfeomorfizmem.
\end{przyklad}

\noindent\textbf{Granica tej części.}
Ciąg~\eqref{eq:exact-surgery-23} kończy
program \emph{wysokowymiarowej} chirurgii
w sensie logicznym: podaje kryterium
istnienia struktury i opis niejednoznaczności.
Nie jest automatyczną tabelą dyfeomorfizmów.
Do obliczenia konkretnego $\mathcal S^s(X)$
potrzeba znać homotopie $G/O$, grupy $L$
danego pierścienia grupowego i działanie
poprzedniej grupy. W niższych wymiarach
podane hipotezy twierdzeń Walla i
$s$-kobordyzmu nie są spełnione.

\par\smallskip
{\footnotesize\noindent\emph{Dalsza lektura:}
\href{https://webhomes.maths.ed.ac.uk/~v1ranick/books/scm.pdf}
{C.~T.~C. Wall, \emph{Surgery on Compact Manifolds}, §10.2--10.8},
\href{https://www.maths.ed.ac.uk/~v1ranick/surgery/ranicki.pdf}
{A.\ Ranicki, \emph{An Introduction to Algebraic Surgery},
część o ciągu dokładnym chirurgii}.\par}

\clearpage
